% !TeX root = ../main.tex
\chapter{Detalhamento do Projeto}
\label{cap:detalhamento}

\section{Visão geral}

A ferramenta será uma aplicação de linha de comando em C\# para .NET 10. A Figura~\ref{fig:blocos} mostra o fluxo para um caso. Cada proposta parte de um commit identificado. A recuperação é usada na condição C do experimento, enquanto as condições A e B passam diretamente da análise à geração. A validação e a emissão do relatório são comuns às três condições.

\begin{figure}[H]
\centering
\begin{tikzpicture}[
    node distance=5mm and 12mm,
    caixa/.style={draw, rounded corners=2pt, align=center, text width=6.2cm, minimum height=8mm, font=\small},
    entrada/.style={caixa, fill=gray!12},
    seta/.style={-{Stealth[length=2mm]}, thick}
]
\node[entrada] (entrada) {Repositório e alerta SARIF ou manifesto};
\node[caixa, below=of entrada] (analysis) {Analysis\\Localização e inventário de autorização};
\node[caixa, below=of analysis] (retrieval) {Retrieval\\Seleção de evidências do repositório};
\node[caixa, below=of retrieval] (generation) {Generation\\Proposta de correção pelo LLM};
\node[caixa, below=of generation] (validation) {Validation\\Compilação, testes e aderência};
\node[caixa, below=of validation] (reporting) {Reporting\\Relatório e revisão humana};
\draw[seta] (entrada) -- (analysis);
\draw[seta] (analysis) -- node[right,font=\footnotesize]{C} (retrieval);
\draw[seta] (retrieval) -- (generation);
\draw[seta] (analysis.east) -- ++(10mm,0) |- node[pos=.29,right,font=\footnotesize]{A e B} (generation.east);
\draw[seta] (generation) -- (validation);
\draw[seta] (validation) -- (reporting);
\end{tikzpicture}
\caption{Fluxo do protótipo de remediação.}
\label{fig:blocos}
\par\smallskip{\footnotesize Fonte elaborada pelos autores.}
\end{figure}

\section{Interfaces entre os blocos}

\begin{description}
    \item[Analysis] Recebe o repositório, o commit e o SARIF ou manifesto. Produz um alerta normalizado, a localização do endpoint e um inventário de autorização.
    \item[Retrieval] Recebe o endpoint e o inventário. Retorna trechos com arquivo, intervalo de linhas, símbolo e justificativa de relevância.
    \item[Generation] Recebe o alerta e o código permitido pela condição, incluindo as evidências recuperadas quando aplicável. Retorna \textit{diff}, justificativa, referências, incerteza e dados de uso do modelo.
    \item[Validation] Recebe a proposta, o commit e a configuração de avaliação. Retorna o estado da aplicação do \textit{diff}, da compilação e dos testes, além dos itens para revisão arquitetural.
    \item[Reporting] Reúne as entradas, a proposta e os resultados de validação. Emite o registro da execução e o material para revisão humana.
\end{description}

\section{Análise do código}

A entrada principal será SARIF 2.1.0 \cite{githubsarif}. O importador extrairá identificador da regra, mensagem, arquivo e região. Uma tabela de correspondência, configurada para cada fonte, associará regras conhecidas a CWE-862 ou CWE-863. Quando a cobertura do analisador for insuficiente, um manifesto experimental indicará identificador do caso, commit, CWE, arquivo e intervalo do endpoint. A origem da entrada permanecerá registrada no relatório. Localização ambígua, classe desconhecida ou caminho externo ao repositório interrompem a geração e produzem um diagnóstico.

Em ambiente isolado, com SDK e dependências preparados, o Roslyn carregará o projeto, localizará a região informada e consultará árvores sintáticas e modelos semânticos \cite{roslyn}. Para controladores, serão registrados classe, ação, rota e método HTTP. Para Minimal APIs, serão registrados o mapeamento, o grupo de rotas, quando resolvível, e a função ou expressão que implementa o endpoint.

O inventário reunirá os metadados de \texttt{Authorize}, \texttt{AllowAnonymous} e \texttt{RequireAuthorization}, os registros feitos por \texttt{AddAuthorization}, as políticas padrão e de \textit{fallback}, papéis, \textit{claims}, implementações de \texttt{IAuthorizationRequirement} e \texttt{IAuthorizationHandler} e chamadas a \texttt{AuthorizeAsync}. Os vínculos serão associados ao nível em que atuam, como configuração global, grupo, controlador, ação ou serviço. Registros condicionais, políticas dinâmicas e dependências que não puderem ser resolvidas serão sinalizados para revisão. A presença de um atributo isolado é insuficiente para classificar a autorização efetiva de um endpoint.

\section{Recuperação contextual}

O recuperador usa o inventário para partir do endpoint e seguir referências a símbolos que definem suas políticas, seus requisitos, seus \textit{handlers} e os serviços de autorização. Uma busca textual complementa os vínculos estruturais para localizar endpoints que operam sobre o mesmo tipo de recurso ou aplicam verificações comparáveis. Os candidatos são ordenados por ligação semântica direta, equivalência da operação e, por último, correspondência textual. O bloco elimina duplicatas e respeita um orçamento de contexto fixado no piloto. Para cada trecho, registra a origem e o motivo da seleção.

Um padrão usado em outro endpoint será considerado evidência somente quando a operação e as condições de acesso forem compatíveis. Por exemplo, editar um documento pode exigir um requisito associado à entidade carregada, enquanto listar documentos pode depender de outra regra. Trechos contraditórios ou insuficientes serão identificados como tal. O primeiro protótipo usará recuperação estrutural e textual no repositório local.

\section{Geração da proposta}

A condição experimental determina se a entrada contém o método ou manipulador associado, o arquivo completo ou o arquivo completo acrescido dos trechos recuperados. As três condições compartilham o mesmo modelo, os parâmetros de geração, as instruções e o formato de saída. O adaptador de LLM será selecionado por configuração. Antes do envio, a CLI limita o conteúdo aos arquivos autorizados e verifica a presença de segredos e configurações sensíveis. Casos que exigirem conteúdo indisponível serão sinalizados.

O modelo deve produzir um \textit{diff} unificado, explicar qual regra foi aplicada e apontar as evidências disponíveis. Quando não houver evidência suficiente para sustentar uma política ou condição de negócio, a saída prevista é uma indicação de incerteza, com possibilidade de abstenção. Os testes do avaliador e a correção de referência ficam fora da entrada. A saída é interpretada como dados e passa por validação de formato antes de qualquer tentativa de aplicação.

\section{Validação e relatório}

A validação cria um Git \textit{worktree} temporário no commit do caso. O \textit{diff} deve ser aplicável e limitado aos arquivos de produção autorizados. Testes e configurações ficam protegidos. Em contêiner sem rede, com dependências preparadas, recursos limitados e sem credenciais, o validador executa \texttt{dotnet build} e os testes funcionais e de autorização. Registra o resultado, a duração e o motivo de falha de cada etapa e descarta a cópia de trabalho.

A revisão de aderência confronta o \textit{diff}, as evidências e a regra esperada, conforme a Seção~\ref{sec:criterios}. O relatório JSON registra alerta, commit, condição, localização, contexto, proposta, estado de validação, tempos e dados do modelo, incluindo custo quando disponível. A versão em Markdown apoia a revisão humana. A decisão de aplicar a correção cabe ao desenvolvedor.

\section{Conjunto de casos e piloto}

Cada caso registra origem, commit, endpoint, CWE, evidências relevantes, regra de acesso, testes de segurança e funcionais e correção de referência. O manifesto de localização fica separado do material de avaliação. A correção de referência documenta uma solução válida, e os testes admitem outros \textit{diffs} que satisfaçam os critérios.

O primeiro caso usará uma API de teste para documentos. Uma operação de alteração depende da regra de propriedade já aplicada em um endpoint equivalente por \texttt{IAuthorizationService} e um \textit{handler}. A avaliação verifica se a proposta reutiliza esse mecanismo, preserva o acesso autorizado e impede efeitos indevidos quando o acesso é negado.

O piloto reunirá cinco a dez casos em duas a três aplicações públicas ou de teste, cobrindo ausência, política incorreta e autorização por recurso. A ampliação para 30 a 50 casos depende de projetos reproduzíveis, regras de acesso documentadas e entradas SARIF ou manifestos.
